\section{Tartışma}
Nicel sonuçlar, modelin test kümesinde dengeli bir precision/recall profili sergilediğini göstermektedir (Tablo~\ref{tab:test_macro}). Bununla birlikte, micro metriklerin macro metriklerden kısmen yüksek olması (Tablo~\ref{tab:test_micro}), daha yoğun damar içeren veya daha büyük alan kaplayan görüntülerin toplam piksel hesabında daha baskın olabildiğine işaret eder. Bu nedenle iki raporlama biçimi birbirini tamamlayıcı olarak değerlendirilmelidir.

Nitel örneklerde, hataların önemli bir kısmının ince damar uçlarında ve düşük kontrast bölgelerde yoğunlaştığı gözlenmiştir (Şekil~\ref{fig:qualitative}). Bu durum, segmentasyon metrikleri kabul edilebilir seviyede olsa dahi, damar topolojisinin tam korunmasını zorlaştırabilir. Nitekim damar analizi sonuçları, damar yoğunluğunun ortalamada benzer kaldığını, ancak iskelet uzunluğu ve dallanma sayısında azalma eğilimi olduğunu göstermektedir (Tablo~\ref{tab:vessel_analysis}). Bu tablo, özellikle ince damarların kopması veya süreksizleşmesi durumunda topolojik ölçümlerin daha duyarlı bir gösterge olabileceğini ortaya koyar.

Öte yandan, test kümesinde iki örnekte TP=FP=FN=0 durumu gözlenmiştir. Bu gibi “tam boş maske” örneklerinde metrikler tanım gereği 1.0 olarak hesaplanabilmektedir. Bu durum veri seti etiketleme içeriği veya ilgili örneklerde damar anotasyonunun bulunmaması gibi nedenlerle ortaya çıkabilir; raporlamada bu özel durum açıkça belirtilmelidir.

Bu çalışmanın başlıca sınırlılıkları; tek bir mimari (U-Net) ile sınırlı kalınması, sabit eşikleme stratejisinin (0.5) tüm görüntüler için optimal olmayabilmesi, ayrıca çıkarılan morfolojik göstergelerin daha ileri post-process (ör. küçük bileşen temizleme, skeleton pruning) adımları ile iyileştirilebilecek olmasıdır. Buna rağmen çalışma, FOV tabanlı ön işleme ve segmentasyon sonrası analiz ile uçtan uca tekrarlanabilir bir damar analizi hattı sunmaktadır.
